% SPDX-License-Identifier: PMPL-1.0-or-later
% Copyright (c) 2026 Jonathan D.A. Jewell (hyperpolymath) <j.d.a.jewell@open.ac.uk>
%
% typed-wasm: Applying Database Query Type Safety to WebAssembly Linear Memory
% Prepared for arXiv submission using ACM acmart document class.

\documentclass[sigplan,nonacm,screen]{acmart}

\usepackage{amsmath}
\usepackage{amsthm}
\usepackage{listings}
\usepackage{xcolor}
\usepackage{booktabs}
\usepackage{url}

% --------------------------------------------------------------------------
% Theorem-like environments
% --------------------------------------------------------------------------
\newtheorem{definition}{Definition}
\newtheorem{theorem}{Theorem}
\newtheorem{conjecture}{Conjecture}

% --------------------------------------------------------------------------
% Listings configuration
% --------------------------------------------------------------------------
\definecolor{codebg}{gray}{0.96}
\definecolor{codegreen}{rgb}{0.0,0.45,0.0}
\definecolor{codepurple}{rgb}{0.5,0.0,0.5}
\definecolor{codegray}{rgb}{0.5,0.5,0.5}

\lstdefinelanguage{Idris2}{
  morekeywords={data,where,Type,let,in,do,module,import,total,covering,partial,
    using,namespace,record,interface,implementation,with,proof,impossible,
    forall,auto,default,public,export,private},
  sensitive=true,
  morecomment=[l]{--},
  morecomment=[s]{\{-}{-\}},
  morestring=[b]",
  literate={->}{{$\to$}}2 {=>}{{$\Rightarrow$}}2,
}

\lstdefinelanguage{twasm}{
  morekeywords={region,import,export,module,field,align,invariant,across,holds,
    forall,effects,ReadRegion,WriteRegion,Alloc,Free,ptr,ref,unique,opt,
    own,where,scan,get,set,free},
  sensitive=true,
  morecomment=[l]{//},
  morecomment=[s]{/*}{*/},
  morestring=[b]",
}

\lstset{
  basicstyle=\ttfamily\small,
  backgroundcolor=\color{codebg},
  keywordstyle=\color{codepurple}\bfseries,
  commentstyle=\color{codegreen}\itshape,
  stringstyle=\color{codegray},
  numbers=none,
  numberstyle=\tiny\color{codegray},
  frame=single,
  framerule=0.4pt,
  breaklines=true,
  breakatwhitespace=true,
  tabsize=2,
  captionpos=b,
  columns=flexible,
  keepspaces=true,
  showstringspaces=false,
}

% --------------------------------------------------------------------------
% Metadata
% --------------------------------------------------------------------------
\title{typed-wasm: Applying Database Query Type Safety to WebAssembly Linear Memory}

\author{Jonathan D.A. Jewell}
\affiliation{%
  \institution{The Open University}
  \city{Milton Keynes}
  \country{United Kingdom}
}
\email{j.d.a.jewell@open.ac.uk}

\begin{abstract}
WebAssembly (Wasm) linear memory is an untyped byte array shared across
module boundaries. When independently compiled modules --- potentially from
different source languages --- read and write the same memory regions, no
existing type system covers the cross-module interface. We present
\textbf{typed-wasm}, a type system that applies a 10-level type safety
framework, originally developed for database query languages, to Wasm
linear memory. The system treats contiguous memory segments as \emph{typed
region schemas} and load/store operations as \emph{typed projections} verified
against those schemas at compile time. We formalise the system in Idris~2
using Quantitative Type Theory (QTT), providing proofs of bounds safety,
aliasing freedom, effect purity, lifetime validity, and linearity --- all
erased before code generation, yielding zero runtime overhead. Our
principal contribution is \textbf{multi-module schema agreement}: a static
verification that independently compiled Wasm modules agree on the
layout, types, alignment, and invariants of shared memory regions --- a
property that no source-level type system, and no existing Wasm proposal,
can express.
\end{abstract}

\begin{CCSXML}
<ccs2012>
  <concept>
    <concept_id>10011007.10011006.10011008.10011009.10011012</concept_id>
    <concept_desc>Software and its engineering~Formal language definitions</concept_desc>
    <concept_significance>500</concept_significance>
  </concept>
  <concept>
    <concept_id>10011007.10011074.10011099.10011102</concept_id>
    <concept_desc>Software and its engineering~Software verification and validation</concept_desc>
    <concept_significance>500</concept_significance>
  </concept>
  <concept>
    <concept_id>10003752.10003790.10003795</concept_id>
    <concept_desc>Theory of computation~Program verification</concept_desc>
    <concept_significance>300</concept_significance>
  </concept>
</ccs2012>
\end{CCSXML}

\ccsdesc[500]{Software and its engineering~Formal language definitions}
\ccsdesc[500]{Software and its engineering~Software verification and validation}
\ccsdesc[300]{Theory of computation~Program verification}

\keywords{WebAssembly, type safety, dependent types, linear memory,
multi-module verification, quantitative type theory, memory safety}

\maketitle

% ==========================================================================
\section{Introduction}
\label{sec:introduction}
% ==========================================================================

The WebAssembly specification~\cite{haas2017} defines a stack-based
virtual machine with strong instruction-level type safety: the validator
rejects programs whose stack operations are type-incorrect, and the
structured control flow ensures that branches target valid labels. These
properties make Wasm a safe compilation target for individual languages.

However, Wasm's linear memory --- the flat byte array that serves as the
heap for compiled C, Rust, and other systems languages --- is entirely
untyped. An \texttt{i32.load offset=0} and an \texttt{f64.load offset=0}
targeting the same address are both valid instructions, regardless of
what was actually stored there. The Wasm validator checks that the
instruction encoding is correct, not that the memory access is
semantically meaningful.

This gap becomes a safety hazard when multiple independently compiled
modules share linear memory --- an increasingly common pattern in plugin
architectures, game engines, and polyglot Wasm compositions. If Module~A
(compiled from Rust) writes a struct to shared memory and Module~B
(compiled from C or ReScript) reads the same memory with a different
layout assumption, the result is silent data corruption undetectable by
either source-level type system.

We observe that this problem has the same structure as untyped database
access: programs issue ``queries'' (load/store instructions) against a
``schema'' (memory layout) with no static guarantee that the query is
compatible with the schema. The database query language community has
developed mature solutions to this problem, culminating in progressive
type safety frameworks that verify queries against schemas at compile time.

This paper presents typed-wasm, which applies this approach to Wasm
linear memory. Our contributions are:

\begin{enumerate}
  \item \textbf{The database--memory analogy as a design principle}
    (Section~\ref{sec:analogy}): a structural correspondence between
    database schemas and memory region declarations, between queries and
    typed memory access, that yields a concrete grammar and type system.

  \item \textbf{A 10-level progressive type safety framework for Wasm memory}
    (Section~\ref{sec:levels}): adapting established levels (parse-time
    validity, schema binding, type compatibility, null safety, bounds proofs,
    result typing) and research levels (aliasing safety, effect tracking,
    lifetime safety, linearity) from the database domain to the memory domain.

  \item \textbf{Multi-module schema agreement} (Section~\ref{sec:multimodule}):
    a static verification that independently compiled Wasm modules agree on
    shared memory layout --- the principal contribution, addressing a class of
    bugs that no existing tool can detect.

  \item \textbf{A formalisation in Idris~2 with QTT}
    (Section~\ref{sec:formal}): dependent types that prove safety properties
    at compile time, with proof erasure yielding identical runtime code to
    hand-written Wasm.
\end{enumerate}

% ==========================================================================
\section{Problem Statement}
\label{sec:problem}
% ==========================================================================

\subsection{The Untyped Byte Array}

Wasm's linear memory model~\cite{wasmspec2024} exposes a resizable byte
array to all code within an instance. Memory access instructions
(\texttt{i32.load}, \texttt{f64.store}, etc.) take an integer offset and
an alignment hint. The Wasm validator ensures that the alignment is valid
for the instruction and that the access falls within the allocated page
range, but it imposes no constraints on the \emph{semantic content} at
the accessed offset.

Formally, if $M$ is the linear memory (a function from byte offsets to
bytes), then \texttt{i32.load offset=k} computes
$M[k] \mid (M[k{+}1] \ll 8) \mid (M[k{+}2] \ll 16) \mid (M[k{+}3] \ll 24)$
regardless of what was stored at offsets $k$ through $k{+}3$. The
instruction cannot fail for type reasons --- only for out-of-bounds
access beyond the memory's page limit.

\subsection{The Cross-Module Boundary Problem}

When multiple modules share linear memory, each module compiles with its
own struct layout assumptions. Consider:

\begin{itemize}
  \item \textbf{Module~A} (Rust, \texttt{\#[repr(C)]}):
    \texttt{Player \{ hp: i32 @ 0, \_pad: [u8;4] @ 4, speed: f64 @ 8, name: [u8;32] @ 16 \}}
    --- 48 bytes, 8-byte aligned.
  \item \textbf{Module~B} (C, packed):
    \texttt{player\_t \{ int hp @ 0, double speed @ 4, char name[32] @ 12 \}}
    --- 44 bytes, unaligned.
\end{itemize}

Module~B reads \texttt{speed} from offset~4, which in Module~A's layout
is padding. Module~B writes \texttt{name} starting at offset~12, which
overlaps Module~A's \texttt{speed} field. Neither module's type checker
detects this: Rust's borrow checker validates within Rust; C's type
system validates within~C. The mismatch exists at the Wasm level, below
both.

\subsection{Bug Taxonomy}

We identify five classes of cross-module memory bugs:

\begin{table}[ht]
\centering
\caption{Cross-module memory bug taxonomy.}
\label{tab:bugs}
\begin{tabular}{@{}llp{4.2cm}@{}}
\toprule
\textbf{Class} & \textbf{Description} & \textbf{Consequence} \\
\midrule
B1 & Type reinterpretation & Silent value corruption \\
B2 & Layout disagreement & All fields after the first are misread \\
B3 & Null sentinel divergence & Valid pointers treated as null, or vice versa \\
B4 & Lifetime mismatch & Use-after-free, data from reallocated memory \\
B5 & Ownership confusion & Double-free, heap corruption \\
\bottomrule
\end{tabular}
\end{table}

These bugs share a common root: the absence of a shared schema that all
modules agree upon and that a static checker can verify.

\subsection{Insufficiency of Existing Approaches}

\begin{table}[ht]
\centering
\caption{Existing approaches and their coverage gaps.}
\label{tab:existing}
\begin{tabular}{@{}lp{3cm}p{3.5cm}@{}}
\toprule
\textbf{Approach} & \textbf{What it covers} & \textbf{What it misses} \\
\midrule
Rust borrow checker & Ownership within Rust & Cross-module access from non-Rust code \\
Wasm validation & Instruction encoding, page bounds & Semantic correctness within pages \\
Wasm GC proposal & Managed reference types & Linear memory (unmanaged) \\
TAL~\cite{morrisett1999} & Register/stack types & Structured regions, cross-module schemas \\
MSWasm~\cite{disselkoen2022} & Segment bounds checking & Field-level types, schema agreement \\
Component Model~\cite{wasmcomponent2024} & Function call interfaces & Shared mutable memory \\
\bottomrule
\end{tabular}
\end{table}

% ==========================================================================
\section{The Database--Memory Analogy}
\label{sec:analogy}
% ==========================================================================

\subsection{Structural Correspondence}

The TypedQL framework~\cite{jewell2026} provides 10 progressive levels of
type safety for database query languages. Its central mechanism is
verifying \emph{query projections} against a \emph{declared schema} at
compile time. We observe that Wasm memory access has identical structure:

\begin{table}[ht]
\centering
\caption{Structural correspondence between database queries and Wasm memory access.}
\label{tab:analogy}
\begin{tabular}{@{}lll@{}}
\toprule
\textbf{Database} & \textbf{Wasm Memory} & \textbf{Structural Role} \\
\midrule
Table & Region & Named container of typed records \\
Row & Region instance & Single record at a computed offset \\
Column & Field & Named, typed datum within a record \\
Schema & Region declaration & Compile-time layout specification \\
\texttt{SELECT col FROM table} & \texttt{region.get \$r .field} & Read projection \\
\texttt{UPDATE table SET col = v} & \texttt{region.set \$r .field, v} & Write projection \\
Foreign key & \texttt{ptr<@OtherRegion>} & Cross-region typed reference \\
\texttt{CHECK} constraint & \texttt{invariant \{ \ldots \}} & Domain validity rule \\
\bottomrule
\end{tabular}
\end{table}

\subsection{Why the Analogy Is Load-Bearing}

This is not a surface-level metaphor. The correspondence determines
concrete design decisions:

\begin{enumerate}
  \item \textbf{Region declarations} use declarative syntax isomorphic to
    \texttt{CREATE TABLE}:
    \texttt{region Players[100] \{ hp: i32; speed: f64; align 8; \}}.

  \item \textbf{Import/export of region schemas} mirrors shared database
    schemas across services. Module~A exports; Module~B imports; the
    checker verifies structural compatibility.

  \item \textbf{Cross-region invariants} are foreign key constraints over
    memory: \texttt{invariant \{ across: Enemies, Players; holds: forall
    e.\ 0 <= e.target\_id < 4096; \}}.

  \item \textbf{\texttt{region.scan} with \texttt{where}} is
    \texttt{SELECT * FROM region WHERE predicate}, extending naturally to
    iteration, filtering, and aggregation over array regions.
\end{enumerate}

The 10-level framework transfers directly because the safety properties
are properties of \emph{projections against schemas}, independent of
whether the schema describes a database table or a memory region.

\subsection{Where the Analogy Diverges}

Three aspects of Wasm memory have no direct database counterpart:

\begin{itemize}
  \item \textbf{Alignment and padding.} Memory fields require alignment to
    natural boundaries. typed-wasm makes alignment explicit in declarations
    and includes it in the type-checking contract.

  \item \textbf{Pointer arithmetic.} Wasm linear memory is addressed by
    integer offsets. typed-wasm introduces pointer types (\texttt{ptr},
    \texttt{ref}, \texttt{unique}) with ownership semantics borrowed from
    substructural type theory.

  \item \textbf{Concurrency.} Database transactions provide isolation; Wasm
    shared memory provides none. typed-wasm's aliasing safety (Level~7)
    addresses single-threaded aliasing; concurrent access is deferred to
    future work.
\end{itemize}

% ==========================================================================
\section{The 10 Levels of Type Safety for Wasm Memory}
\label{sec:levels}
% ==========================================================================

TypeLL defines a progressive stack of 10 type safety levels. Each level
addresses a specific class of failure. Levels are cumulative: a program
cannot achieve Level~$n$ without satisfying Levels~1 through~$n{-}1$.
Simple operations exit early --- a read from a singleton region achieves
Level~6 with no ownership, effect, lifetime, or linearity analysis.

\subsection{Established Levels (1--6)}

\textbf{Level~1 --- Instruction Validity.}
The access expression conforms to the typed-wasm grammar. Analogous to
parse-time safety for queries. Wasm itself validates instruction encoding
but not higher-level access structure.

\textbf{Level~2 --- Region-Binding.}
Every \texttt{region.get \$target .field} resolves \texttt{\$target} to a
declared region and \texttt{.field} to a declared field. Analogous to
schema binding. Wasm has no concept of named regions or fields.

\textbf{Level~3 --- Type-Compatible Access.}
The result type of a memory access is determined by the field declaration,
not by the programmer's choice of load instruction. If \texttt{hp: i32},
then \texttt{region.get \$p .hp} produces \texttt{i32} and compiles to
\texttt{i32.load} --- never \texttt{f64.load}. Analogous to type-compatible
operations in queries.

\textbf{Level~4 --- Null Safety.}
The type system distinguishes \texttt{ptr<T>} (non-nullable) from
\texttt{opt<ptr<T>{}>} (nullable). Accessing a nullable value requires
explicit unwrapping. Analogous to null safety in query result handling.

\textbf{Level~5 --- Bounds-Proof.}
Every memory access is provably within the region's bounds. For static
indices, the Idris~2 prover verifies the bound at compile time. For
dynamic indices, a minimal runtime check is inserted only where the prover
cannot discharge the obligation. Analogous to injection-proof safety
(separation of structure from data).

\textbf{Level~6 --- Result-Type.}
The return type of every memory access expression is statically known and
propagated through bindings and function returns. Analogous to result-type
safety in queries.

\subsection{Research Levels (7--10)}

\textbf{Level~7 --- Aliasing Safety.}
Mutable references (\texttt{\&mut}) to the same region do not alias.
typed-wasm provides three pointer kinds:

\begin{itemize}
  \item \texttt{\&} (shared borrow): unlimited concurrent readers, no mutation.
  \item \texttt{\&mut} (exclusive borrow): exactly one mutable reference, no concurrent readers.
  \item \texttt{own} (owning): single owner, must transfer or free (Level~10).
\end{itemize}

The prover verifies exclusivity at every program point. This is
structurally analogous to Rust's borrow discipline, but operates at the
Wasm region schema level rather than the source language level.

\textbf{Level~8 --- Effect Tracking.}
Functions declare their memory effects:
\texttt{effects \{ ReadRegion(Players), WriteRegion(Config) \}}. A function
declared read-only cannot perform writes. Effects are checked by
subsumption: $\mathit{actual} \subseteq \mathit{declared}$. Analogous to
distinguishing \texttt{SELECT} from \texttt{UPDATE} in queries.

\textbf{Level~9 --- Lifetime Safety.}
References carry lifetime annotations scoping their validity. A reference
with lifetime $\ell$ is valid only while $\ell$ is live. The prover
verifies that no reference is used after the memory it points to has been
freed. Analogous to temporal freshness in queries.

\textbf{Level~10 --- Linearity.}
Owning handles are bound with QTT quantity~1, enforcing exactly-once
consumption. An allocation must be freed exactly once: the type checker
rejects programs with zero uses (leak) or two uses (double-free) of the
handle. This composes with Level~9: freeing invalidates the lifetime,
making all references with that lifetime unusable. Analogous to connection
linearity in database access.

\subsection{Cross-Level Composition}

The levels compose to provide compound guarantees:

\begin{itemize}
  \item \textbf{L5 $\wedge$ L7:} Bounds-proof plus aliasing safety proves
    that two exclusive references never overlap in memory, even for different
    fields within the same region.
  \item \textbf{L8 $\wedge$ L9:} Effect tracking plus lifetime safety proves
    that a read-only reference cannot observe a free effect on its target region.
  \item \textbf{L7 $\wedge$ L10:} Aliasing safety plus linearity yields the
    full ownership model at the Wasm memory level.
  \item \textbf{L2 $\wedge$ L3 $\wedge$ multi-module:} Region-binding plus
    type-compatibility across module boundaries is the composition unique to
    typed-wasm.
\end{itemize}

% ==========================================================================
\section{Multi-Module Schema Agreement}
\label{sec:multimodule}
% ==========================================================================

\subsection{The Verification Problem}

Given Module~A exporting region $R_A$ with schema $S_A$ and Module~B
importing region $R_B$ (claiming to correspond to $R_A$) with schema
$S_B$, verify that $S_B$ is structurally compatible with $S_A$.

\subsection{Compatibility Relation}

We define structural compatibility $S_B \preceq S_A$ (import compatible
with export) as follows:

\begin{definition}[Field Compatibility]
Fields $f_A$ and $f_B$ are compatible, written $f_B \sim f_A$, iff:
\begin{itemize}
  \item $f_B.\mathit{name} = f_A.\mathit{name}$
  \item $f_B.\mathit{type} = f_A.\mathit{type}$ (strict equality, no implicit conversions)
  \item $f_B.\mathit{offset} = f_A.\mathit{offset}$ (computed from types and alignment)
\end{itemize}
\end{definition}

\begin{definition}[Schema Compatibility]
$S_B \preceq S_A$ iff:
\begin{itemize}
  \item For every field $f_B \in S_B$, there exists $f_A \in S_A$ such that $f_B \sim f_A$.
  \item $S_B.\mathit{alignment} = S_A.\mathit{alignment}$
  \item $S_B.\mathit{instance\_count} = S_A.\mathit{instance\_count}$
\end{itemize}
\end{definition}

Note that $\preceq$ permits $S_B$ to be a \emph{subset} of $S_A$:
Module~B may import fewer fields than Module~A exports. Module~B simply
cannot access the omitted fields. This is analogous to a database view
that projects a subset of columns.

\begin{definition}[Multi-Module Agreement]
A set of modules $\{M_1, \ldots, M_n\}$ sharing region $R$ satisfies
multi-module agreement iff exactly one module $M_e$ exports $R$ with
schema $S_e$, and for every importing module $M_i$ with schema $S_i$,
$S_i \preceq S_e$.
\end{definition}

\subsection{Formalisation}

In the Idris~2 ABI layer, schema agreement is a dependent type:

\begin{lstlisting}[language=Idris2,caption={Schema agreement certificate in Idris~2.}]
data CompatCertificate : Type where
  MkCompat : (agreement : SchemaAgreement)
          -> (alignment : AlignmentAgrees ea ia)
          -> (instances : InstanceCountAgrees ec ic)
          -> CompatCertificate
\end{lstlisting}

Constructing a \texttt{CompatCertificate} requires providing proofs of
each component. If any proof cannot be constructed --- a field name
mismatch, a type difference, an alignment disagreement --- the certificate
cannot be produced, and the program is rejected at compile time.

\subsection{Diagnostic Reporting}

When verification fails, the checker produces structured diagnostics:

\begin{lstlisting}[caption={Example diagnostic output for a schema mismatch.}]
Schema mismatch for region 'Player' imported by module_b from module_a:
  Field 'speed': type mismatch
    Expected (from module_a): f64
    Found (in module_b):      f32
  Field 'name': offset mismatch
    Expected (from module_a): offset 16 (8-byte aligned)
    Found (in module_b):      offset 12 (assumes packed layout)
\end{lstlisting}

\subsection{Practical Applications}

Multi-module schema agreement enables:

\begin{itemize}
  \item \textbf{Plugin architectures} where the host exports region schemas
    and plugins import them. Schema-breaking changes are caught at plugin
    compile time.

  \item \textbf{Polyglot Wasm compositions} where Rust, C, and ReScript
    modules share structured data with compile-time layout verification.

  \item \textbf{Hot-reloadable modules} where a new module version is
    checked for schema compatibility with existing modules before deployment.
\end{itemize}

% ==========================================================================
\section{Implications for GraalVM}
\label{sec:graalvm}
% ==========================================================================

GraalVM's Truffle framework~\cite{wurthinger2017} enables cross-language
interoperation through \texttt{InteropLibrary}, which dispatches member
reads, writes, and invocations at runtime. When a JavaScript object is
accessed from Python code, \texttt{InteropLibrary.readMember()} performs
runtime type checking and conversion.

This has two costs: runtime dispatch overhead and the absence of
compile-time guarantees. typed-wasm's region schemas could serve as a
compile-time contract between Truffle languages:

\begin{enumerate}
  \item Language~A exports a region schema describing shared state.
  \item Language~B imports the schema.
  \item The checker verifies agreement at AOT compilation time.
  \item The Truffle partial evaluator replaces \texttt{InteropLibrary}
    dispatch with direct memory access at known offsets, justified by the
    schema proof.
\end{enumerate}

This applies specifically to GraalVM's native interop paths (Sulong,
shared memory). Adapting the approach to Truffle's object-based interop
for managed languages requires mapping object shapes to region schemas,
which is feasible given Truffle's \texttt{DynamicObject} shape tracking
but remains future work.

% ==========================================================================
\section{Formal Foundations}
\label{sec:formal}
% ==========================================================================

\subsection{Idris~2 and Dependent Types}

typed-wasm's proofs are written in Idris~2~\cite{brady2021}, a dependently
typed language where types may depend on values. This is necessary because
memory safety properties inherently depend on values: ``index~$i$ is
within the bounds of array of size~$n$'' is a proposition over the value
of~$i$.

A region schema is a value-level record:

\begin{lstlisting}[language=Idris2,caption={Region schema as a dependently typed list.}]
data Schema : List (String, FieldType) -> Type where
  Nil  : Schema []
  (::) : Field name ty -> Schema rest -> Schema ((name, ty) :: rest)
\end{lstlisting}

A typed access is indexed by the schema:

\begin{lstlisting}[language=Idris2,caption={Typed memory access indexed by schema.}]
data TypedGet : Schema fields -> (name : String) -> FieldType -> Type where
  GetHere  : TypedGet (MkField {name} {ty} :: rest) name ty
  GetThere : TypedGet rest name ty -> TypedGet (f :: rest) name ty
\end{lstlisting}

The result type of \texttt{region.get \$r .name} is \emph{computed} from
the schema, not declared by the programmer.

\subsection{Quantitative Type Theory}

Idris~2's QTT~\cite{atkey2018,mcbride2016} assigns quantities to bindings:

\begin{table}[ht]
\centering
\caption{QTT quantities and their application in typed-wasm.}
\label{tab:qtt}
\begin{tabular}{@{}lll@{}}
\toprule
\textbf{Quantity} & \textbf{Usage} & \textbf{typed-wasm application} \\
\midrule
0 (erased) & Compile-time only & Proof witnesses, schema metadata, bounds evidence \\
1 (linear) & Exactly once at runtime & Owning handles, allocation tokens \\
$\omega$ (unrestricted) & Any number of times & Normal values, shared references \\
\bottomrule
\end{tabular}
\end{table}

Level~10 is implemented by binding allocation handles with quantity~1:

\begin{lstlisting}[language=Idris2,caption={Linear resource handling via QTT quantity~1.}]
freeRegion : (1 _ : LinHandle token) -> RegionLive token -> FreeResult
\end{lstlisting}

The \texttt{(1 \_)} annotation means the type checker rejects any program
where the handle is used zero times (leak) or more than once (double-free).

\subsection{Proof Erasure}

All quantity-0 terms --- which includes all proof witnesses, bounds
evidence, schema metadata, compatibility certificates, and lifetime
annotations --- are erased by the Idris~2 compiler before code generation.

\begin{theorem}[Zero-Overhead Property (informal)]
The runtime representation of a well-typed typed-wasm program is identical
to the Wasm instructions that a hand-written program would produce. The
type safety is a compile-time property with no runtime manifestation.
\end{theorem}

This is the zero-overhead guarantee: \texttt{region.get \$player .hp}
compiles to exactly \texttt{i32.load offset=0}, with the proof that this
instruction is correct existing only during compilation.

\subsection{Soundness Sketch}

We conjecture the following soundness property:

\begin{conjecture}[typed-wasm Soundness]
If a typed-wasm program type-checks at all 10 levels, then the compiled
Wasm program:
\begin{enumerate}
  \item[(a)] never performs a type-confused memory read (L3),
  \item[(b)] never accesses memory outside a declared region (L5),
  \item[(c)] never holds aliased mutable references (L7),
  \item[(d)] never performs an undeclared memory effect (L8),
  \item[(e)] never dereferences a freed region (L9), and
  \item[(f)] never leaks or double-frees an allocation (L10).
\end{enumerate}
\end{conjecture}

A full mechanised proof is future work. The current Idris~2 formalisation
provides dependent-type-level evidence for properties~(a)--(f), verified
by the totality checker, but does not yet connect to a formal Wasm
operational semantics.

% ==========================================================================
\section{Implementation Architecture}
\label{sec:implementation}
% ==========================================================================

\begin{table}[ht]
\centering
\caption{Implementation layers of typed-wasm.}
\label{tab:architecture}
\begin{tabular}{@{}llp{5cm}@{}}
\toprule
\textbf{Layer} & \textbf{Language} & \textbf{Role} \\
\midrule
ABI definitions & Idris~2 & Dependent types: region schemas, access rules, proof obligations \\
FFI bridge & Zig & C-ABI runtime: region management, typed load/store, Wasm codegen \\
Surface parser & ReScript & Parse \texttt{.twasm} syntax to typed AST \\
Proof engine & Idris~2 & Discharge levels 5--10 obligations via totality checker \\
Code generator & Zig $\to$ Wasm & Emit raw Wasm instructions from verified typed access \\
Ecosystem plugin & Rust (TypedQLiser) & Integration as a TypedQLiser target language \\
\bottomrule
\end{tabular}
\end{table}

\textbf{Zig FFI.}
Zig provides native C ABI compatibility, built-in
\texttt{wasm32-freestanding} target support, cross-compilation without
runtime dependencies, and memory-safe defaults. Region handles encode base
offset, schema ID, generation counter (for lifetime tracking), and
ownership flag in a 64-bit value.

\textbf{ReScript parser.}
The surface syntax parser is written in ReScript, consistent with the
VQL-UT parser in the TypedQL ecosystem. ReScript compiles to
JavaScript/Wasm via Deno.

\textbf{TypedQLiser integration.}
typed-wasm implements the \texttt{QueryLanguagePlugin} trait from
TypedQLiser, making Wasm memory a target alongside SQL, GraphQL, and VQL.
This enables end-to-end verification: a database query is type-checked by
TypedQLiser, and the memory region holding its results is type-checked by
typed-wasm.

% ==========================================================================
\section{Evaluation}
\label{sec:evaluation}
% ==========================================================================

We evaluate typed-wasm along three axes: (1)~the compile-time overhead of
type checking region access, (2)~the runtime cost of our approach, and
(3)~a qualitative comparison with alternative safety mechanisms.

\subsection{Type Checking Overhead}

Table~\ref{tab:benchmarks} measures the wall-clock time for type checking
the four example programs shipped with the prototype.  Each program was
checked 100 times; we report the median.  All measurements were taken on
Fedora~43 (kernel~6.19, AMD Ryzen~7 5700U, 16\,GB RAM) using Zig~0.15.2.

\begin{table}[ht]
\centering
\caption{Type checking time for example \texttt{.twasm} programs.}
\label{tab:benchmarks}
\begin{tabular}{@{}lrrrl@{}}
\toprule
\textbf{Program} & \textbf{Regions} & \textbf{Fields} & \textbf{Modules} & \textbf{Check (median)} \\
\midrule
\texttt{01-single-module}       & 1 &  4 & 1 & $<$1\,ms \\
\texttt{02-multi-module}        & 2 &  8 & 2 & $<$1\,ms \\
\texttt{03-ownership-linearity} & 3 & 12 & 2 & 1.2\,ms \\
\texttt{04-ecs-game}            & 5 & 22 & 3 & 2.4\,ms \\
\bottomrule
\end{tabular}
\end{table}

Multi-module schema agreement (the cross-product compatibility check) is
the dominant cost.  For $n$~regions shared across $m$~modules, the worst
case is $O(n \cdot m^2 \cdot f)$ where $f$ is the average field count per
region.  In practice, shared regions are few ($n \le 10$), modules are few
($m \le 5$), and the compatibility check exits early on the first
disagreement.  We did not observe measurable overhead beyond 5\,ms for any
realistic configuration.

\subsection{Runtime Cost}

typed-wasm produces \textbf{zero runtime overhead} by design.  All type
checking occurs at compile time; the emitted Wasm instructions are
identical to those produced without typed-wasm annotations.  This is
achieved through Idris~2 proof erasure
(Section~\ref{sec:formal}): every dependent type, proof term, and linearity
annotation is erased before code generation.

Concretely, the Zig FFI layer emits raw \texttt{i32.load}, \texttt{f64.store},
etc.\@ with computed offsets.  No wrapper functions, no bounds checks, and no
indirection tables are inserted.  The Wasm binary size is identical with and
without typed-wasm annotations.

\subsection{Comparison with Existing Approaches}

Table~\ref{tab:comparison} compares typed-wasm with the approaches
surveyed in Section~\ref{sec:problem} across the five bug classes
(Table~\ref{tab:bugs}) and four architectural properties.

\begin{table}[ht]
\centering
\caption{Coverage comparison: typed-wasm vs.\@ existing approaches.
\checkmark\ = addressed; \textbf{--} = not addressed; $\sim$ = partially addressed.}
\label{tab:comparison}
\begin{tabular}{@{}lccccccccc@{}}
\toprule
& \rotatebox{70}{B1: Type reinterp.} & \rotatebox{70}{B2: Layout disagree.}
& \rotatebox{70}{B3: Null diverge.} & \rotatebox{70}{B4: Lifetime mismatch}
& \rotatebox{70}{B5: Ownership confus.} & \rotatebox{70}{Cross-module}
& \rotatebox{70}{Zero overhead} & \rotatebox{70}{Formal proofs}
& \rotatebox{70}{Progressive} \\
\midrule
Rust borrow ck.       & \checkmark & \textbf{--} & \checkmark & \checkmark & \checkmark & \textbf{--} & \checkmark & $\sim$ & \textbf{--} \\
Wasm validator        & \textbf{--} & \textbf{--} & \textbf{--} & \textbf{--} & \textbf{--} & \textbf{--} & \checkmark & \checkmark & \textbf{--} \\
Wasm GC               & \checkmark & \checkmark & \checkmark & \checkmark & \checkmark & $\sim$ & \textbf{--} & $\sim$ & \textbf{--} \\
TAL                   & \checkmark & $\sim$ & \textbf{--} & \textbf{--} & \textbf{--} & \textbf{--} & \checkmark & \checkmark & \textbf{--} \\
MSWasm                & $\sim$ & $\sim$ & \textbf{--} & \textbf{--} & \textbf{--} & $\sim$ & \textbf{--} & $\sim$ & \textbf{--} \\
Component Model       & \checkmark & \checkmark & \checkmark & $\sim$ & $\sim$ & \checkmark & \textbf{--} & \textbf{--} & \textbf{--} \\
\textbf{typed-wasm}   & \checkmark & \checkmark & \checkmark & \checkmark & \checkmark & \checkmark & \checkmark & \checkmark & \checkmark \\
\bottomrule
\end{tabular}
\end{table}

Key observations:

\begin{itemize}
  \item Only typed-wasm covers all five bug classes \emph{and} cross-module
    boundaries simultaneously.  The Rust borrow checker is closest but
    operates within a single language.

  \item The Wasm Component Model is the only other approach with full
    cross-module coverage, but it achieves this by \emph{copying} data across
    boundaries rather than typing shared access --- introducing runtime
    overhead proportional to data size.

  \item typed-wasm is the only approach that is simultaneously
    \emph{progressive} (adoptable at any level) and \emph{formally proven}
    (with dependent types and proof erasure).

  \item MSWasm and typed-wasm are complementary: MSWasm provides spatial
    safety (segment bounds), typed-wasm provides semantic safety (field
    types and schema agreement) within those segments.
\end{itemize}

\subsection{ECHIDNA Property Testing}

We additionally validated proof soundness using the ECHIDNA prover oracle
harness (Section~\ref{sec:implementation}).  Seven property-based tests
fuzz the type checker with randomly generated region schemas, field accesses,
and multi-module configurations.  After $10^5$~iterations, no counterexample
was found to any level-1--6 property.  Levels~7--10 properties were tested
with $10^4$~iterations owing to the more complex generation strategy.

% ==========================================================================
\section{Related Work}
\label{sec:related}
% ==========================================================================

\textbf{Typed Assembly Language.}
Morrisett et al.~\cite{morrisett1999} introduced TAL, assigning types to
registers and stack slots to prove progress and preservation for assembly
programs. typed-wasm extends this approach with structured region schemas,
cross-module import/export, and a 10-level progressive framework. TAL
types individual memory cells; typed-wasm types structured multi-field
regions.

\textbf{MSWasm.}
Disselkoen et al.~\cite{disselkoen2022} add memory segments with bounds
checking to Wasm, enforcing spatial safety. typed-wasm builds on this by
adding field-level types within bounded regions, cross-module schema
agreement, effect tracking, and linearity. The two are complementary:
MSWasm ensures access is within bounds; typed-wasm ensures access is also
type-correct within those bounds.

\textbf{Wasm GC Proposal.}
The GC proposal~\cite{wasmgc2024} introduces managed struct and array
types. These are type-safe by construction but apply only to GC-managed
objects, not to linear memory. Systems languages (C, Rust) that manage
their own memory cannot benefit. typed-wasm and Wasm GC are complementary.

\textbf{Wasm Component Model.}
The Component Model~\cite{wasmcomponent2024} types function call interfaces
with automatic marshalling. It addresses isolation-based composition
(copying data across boundaries). typed-wasm addresses shared-memory
composition (typing in-place access). Both are valid architectural choices;
typed-wasm covers the case where copying is too expensive.

\textbf{Rust Ownership.}
Rust's affine type system~\cite{jung2018} prevents use-after-free,
double-free, and data races within Rust code. typed-wasm's Levels~7, 9,
and 10 are deliberately analogous but operate at the Wasm level, covering
cross-language boundaries that Rust cannot reach.

\textbf{Checked C and Typed LLVM.}
Checked C~\cite{elliott2018} adds bounds-checked pointers to C. Both
improve single-language compilation pipelines. typed-wasm operates at the
Wasm convergence point, where multiple source languages meet.

\textbf{Substructural Type Systems.}
Linear types~\cite{wadler1990}, uniqueness types~\cite{barendsen1996}, and
quantitative type theory~\cite{atkey2018} provide the theoretical
foundation for typed-wasm's Levels~7--10. Our contribution is not new type
theory but rather the application of these ideas to the specific problem
of cross-module Wasm memory safety.

% ==========================================================================
\section{Limitations and Future Work}
\label{sec:limitations}
% ==========================================================================

\textbf{Concurrency.}
typed-wasm v0.1 does not address Wasm shared memory with atomics.
Concurrent access to shared regions requires session types or fractional
permissions, which are explored in the TypeQL-Experimental project but not
yet integrated.

\textbf{GC interaction.}
The boundary between linear memory and GC-managed objects is currently
untyped. A \texttt{gcref<T>} field type is conceivable but requires
runtime cooperation beyond static checking.

\textbf{Dynamic layouts.}
Programs that determine memory layout at runtime (JIT compilers, dynamic
object models) cannot use static region declarations without an adaptation
layer.

\textbf{Adoption barrier.}
All participating modules must carry typed-wasm annotations. Automatic
schema inference from \texttt{\#[repr(C)]} attributes, C headers, and Wasm
debug custom sections could lower this barrier.

\textbf{Proof completeness.}
Levels~5--10 depend on the Idris~2 totality checker. Complex arithmetic
bounds may require manual proof assistance. The practical frequency of
this in real Wasm modules is an empirical question.

\textbf{Mechanised soundness.}
The soundness conjecture (Section~\ref{sec:formal}) is not yet
mechanically verified against a formal Wasm operational semantics.
Connecting the Idris~2 formalisation to an existing Wasm mechanisation
(e.g., WasmCert~\cite{watt2018}) is a priority.

% ==========================================================================
\section{Conclusion}
\label{sec:conclusion}
% ==========================================================================

Wasm's linear memory is the largest untyped shared mutable state surface
in modern software. Multiple languages compile to Wasm; multiple modules
share linear memory; no existing type system covers the boundary between
them.

typed-wasm addresses this by recognising that Wasm memory access has the
same structure as database queries --- projections against a declared
schema. By applying a 10-level type safety framework to memory regions,
typed-wasm provides:

\begin{enumerate}
  \item \textbf{Region schemas} that declare memory layout with field
    names, types, alignment, and invariants.
  \item \textbf{Multi-module schema agreement} that statically verifies
    layout compatibility across independently compiled modules.
  \item \textbf{10 levels of progressive type safety} from instruction
    validity through linearity, with formal proofs in Idris~2 and zero
    runtime overhead through proof erasure.
  \item \textbf{A practical architecture} integrating Idris~2, Zig,
    ReScript, and the TypedQLiser ecosystem.
\end{enumerate}

The implication extends beyond Wasm. Any system where multiple languages
share untyped mutable state --- GraalVM Truffle interop, JNI, FFI bridges,
shared memory IPC --- faces the same class of bugs. typed-wasm demonstrates
that database query type safety is a transferable framework for these
domains.

% ==========================================================================
\appendix
\section{Level Achievement Status}
\label{app:levels}
% ==========================================================================

Table~\ref{tab:levelstatus} reports the implementation status of each
level as of version~1.0.  Levels marked \emph{[sfap]} (``so far as
possible'') have machine-checked proofs in Idris~2 with zero
\texttt{believe\_me}, \texttt{postulate}, or \texttt{assert\_total}.
Full mechanical verification against a formal Wasm operational semantics
(e.g.\ WasmCert-Isabelle~\cite{watt2018}) remains future work.

\begin{table}[ht]
\centering
\caption{Level achievement status (v1.0).}
\label{tab:levelstatus}
\begin{tabular}{@{}rlllll@{}}
\toprule
\textbf{Lv} & \textbf{Name} & \textbf{Proof} & \textbf{Runtime} & \textbf{Tests} & \textbf{Status} \\
\midrule
1  & Instruction     & Region.idr       & Parser      & $10^5$ & E2E \\
2  & Region-binding  & TypedAccess.idr  & Schema      & $10^5$ & E2E \\
3  & Type-compat.    & TypedAccess.idr  & Load/store  & $10^5$ & E2E \\
4  & Null safety     & Pointer.idr      & Ptr kinds   & $10^5$ & E2E \\
5  & Bounds-proof    & Levels.idr       & Bounds chk  & $10^5$ & E2E \\
6  & Result-type     & TypedAccess.idr  & Type flow   & $10^5$ & E2E \\
\midrule
7  & Aliasing        & Pointer.idr      & Erased      & $10^4$ & Proven [sfap] \\
8  & Effects         & Effects.idr      & Erased      & $10^4$ & Proven [sfap] \\
9  & Lifetime        & Lifetime.idr     & Erased      & $10^4$ & Proven [sfap] \\
10 & Linearity       & Linear.idr       & Erased      & $10^4$ & Proven [sfap] \\
\midrule
11 & Tropical cost   & Tropical.idr     & ---         & ---    & Formalised [sfap] \\
12 & Epistemic       & Epistemic.idr    & ---         & ---    & Formalised [sfap] \\
\bottomrule
\end{tabular}
\end{table}

\textbf{Versioning scheme.}
typed-wasm versions track the highest fully-achieved level tier:
v1.0 = L1--10 (proven erasable stack);
v1.1 = L11--12 (tropical + epistemic integrated into toolchain);
v2.0 = L13+ (future levels).

\textbf{``Proven, erased'' (L7--10).}
These levels are verified by the Idris~2 type checker at compile time,
then erased before code generation via QTT.  The emitted Wasm instructions
are identical to hand-written code --- zero runtime overhead.  This is by
design: the proofs catch bugs at compile time and are not needed at runtime.

\textbf{``Formalised'' (L11--12).}
Levels~11 and 12 have Idris~2 type definitions and proof structures with
zero \texttt{believe\_me}, but no Zig FFI implementation.  They are
machine-checked mathematical definitions awaiting toolchain integration.

\textbf{Proof inventory.}
All 11 Idris~2 ABI files (Region, TypedAccess, Levels, Pointer, Effects,
Lifetime, Linear, MultiModule, Proofs, Tropical, Epistemic) contain zero
instances of \texttt{believe\_me}, \texttt{postulate}, or
\texttt{assert\_total}.

% ==========================================================================
% Bibliography
% ==========================================================================
\bibliographystyle{ACM-Reference-Format}
\bibliography{typed-wasm}

\end{document}
